\documentclass[12pt,letterpaper]{article}

\usepackage{amsmath,amssymb,amsthm}
\usepackage{mathtools}
\usepackage{geometry}
\usepackage{hyperref}
\usepackage{booktabs}
\usepackage{array}
\usepackage{xcolor}
\usepackage{framed}
\usepackage{microtype}
\usepackage{enumitem}
\usepackage{mdframed}
\usepackage[numbers,sort&compress]{natbib}

\geometry{margin=1in}

\theoremstyle{plain}
\newtheorem{theorem}{Theorem}[section]
\newtheorem{lemma}[theorem]{Lemma}
\newtheorem{proposition}[theorem]{Proposition}
\theoremstyle{definition}
\newtheorem{definition}[theorem]{Definition}
\newtheorem{openprob}[theorem]{Open Problem}
\theoremstyle{remark}
\newtheorem{remark}[theorem]{Remark}
\newtheorem{flag}[theorem]{Flag}

\definecolor{flagred}{rgb}{0.75,0.1,0.1}
\definecolor{openblue}{rgb}{0.1,0.1,0.75}
\definecolor{boxgray}{rgb}{0.95,0.95,0.95}

\newenvironment{expertflag}{%
  \begin{mdframed}[backgroundcolor=boxgray,linecolor=flagred,linewidth=1.2pt]
  \noindent\textcolor{flagred}{\textbf{Expert Flag.}}\quad}{%
  \end{mdframed}}

\newenvironment{openprobbox}{%
  \begin{mdframed}[backgroundcolor=boxgray,linecolor=openblue,linewidth=1.2pt]
  \noindent\textcolor{openblue}{\textbf{Open Problem.}}\quad}{%
  \end{mdframed}}

\newcommand{\kB}{k_{\mathrm{B}}}
\newcommand{\lP}{\ell_{P}}
\newcommand{\kt}{\tilde{\kappa}}
\newcommand{\kteff}{\tilde{\kappa}_{\mathrm{eff}}}
\newcommand{\Sent}{S_{\mathrm{ent}}}
\newcommand{\Ment}{M_{\mathrm{ent}}}
\newcommand{\asc}{\alpha_{\mathrm{screen}}}
\newcommand{\tr}{\operatorname{Tr}}
\newcommand{\mP}{m_{\mathrm{P}}}
\newcommand{\EP}{E_{\mathrm{P}}}

\title{\textbf{Expert Analysis: Entanglement Entropy Density as a Gravitational Source}\\[6pt]
\large Theoretical Tightening, OP3 Lindblad Model, NISQ Experiment Design,\\
and Referee-Level Critique}

\author{Analysis prepared by: Expert Theoretical Physics Review\\[2pt]
\small Responding to: \textit{Refined Entanglement Gravity White Paper}, K.\ Monette (March 2026)}

\date{March 2026}

\begin{document}
\maketitle
\tableofcontents
\bigskip

% ─────────────────────────────────────────────────────────────────────────────
\section*{Framework Summary (2--3 paragraphs as requested)}
\addcontentsline{toc}{section}{Framework Summary}

The refined white paper proposes adding a stress-energy source term
$T^{\mathrm{ent}}_{\mu\nu} = \tilde{\kappa}(c^4/8\pi G\kB\ln 2)\Sent g_{\mu\nu}$
to the Einstein field equations, where $\Sent$ is the entanglement entropy density of a
quantum system in bit\,m$^{-3}$ and $\kt$ is a dimensionless coupling.  The dimensional
consistency of this addition is established regardless of the value of $\kt$.  The specific
estimate $\kt = -1/4$ is derived under a Planck-scale holographic regulator
hypothesis --- the ansatz that volume-law entanglement entropy in laboratory systems
couples to spacetime geometry through the same Planck-length regulator that governs
horizon area entropy.  This is explicitly labeled a motivated hypothesis, not a theorem.

Three structural results are established correctly.  First, the radial acceleration
correction $\Delta a(R) \propto \Ment(<\!R)/R^2$ for any bounded, spherically symmetric
source --- this resolves the prior confusion about $1/R$ versus $1/R^2$ scaling.  Second,
the Bianchi identity forces $\nabla^\mu T^{\mathrm{matter}}_{\mu\nu} = -\nabla_\nu(\kt c^4 \Sent/8\pi G\kB\ln 2)$,
so matter is not independently conserved in entanglement-entropy gradients; the paper
adopts an ``entanglement-sector exchange'' interpretation where the \emph{total} stress-energy
is conserved.  Third, existing equivalence-principle experiments (MICROSCOPE, Asenbaum et al.)
do not directly constrain $\kt$ because their test masses have negligible and nearly identical
$\Sent$.

The experimental program has been correctly revised: macroscopic $^{87}$Rb BEC interferometry
is demoted because Gaussian spatial screening gives $\asc \sim e^{-(R/\xi)^2} \sim 10^{-43}$
for cloud radius $R \gg$ coherence length $\xi$, making the signal undetectable.  The primary
near-term target is now a 50--100 qubit NISQ platform in the regime $R \lesssim \xi$.
The falsification criterion --- requiring conditions F1--F5 on qubit count, pressure
sensitivity, run count, platform independence, and null result --- is stated without requiring
knowledge of $\asc$ or the Planck-scale renormalization of $\kt$.

% ─────────────────────────────────────────────────────────────────────────────
\section{Tightening the Theoretical Core}
\label{sec:theory}

\subsection{Covariant notation and the full modified field equations}

The white paper's modified field equations are restated here in the most standard
covariant notation.  Greek indices $\mu,\nu,\ldots$ run $0,1,2,3$; signature $(-,+,+,+)$;
$G = 6.674\times10^{-11}$\,m$^3$\,kg$^{-1}$\,s$^{-2}$; $c = 2.998\times10^8$\,m\,s$^{-1}$;
$\kB = 1.381\times10^{-23}$\,J\,K$^{-1}$; $\lP = (\hbar G/c^3)^{1/2} = 1.616\times10^{-35}$\,m.

\textbf{Total stress-energy decomposition.}
\begin{equation}
  T^{\mathrm{total}}_{\mu\nu}
  \;\equiv\;
  T^{\mathrm{matter}}_{\mu\nu}
  \;+\;
  T^{\mathrm{ent}}_{\mu\nu},
  \label{eq:decomp}
\end{equation}
where
\begin{equation}
  T^{\mathrm{ent}}_{\mu\nu}
  \;\equiv\;
  -\Lambda_{\mathrm{ent}}\,g_{\mu\nu},
  \quad
  \Lambda_{\mathrm{ent}}(x)
  \;\equiv\;
  -\,\kt\,\frac{c^4}{8\pi G\kB\ln 2}\,\Sent(x).
  \label{eq:Tent}
\end{equation}
The sign convention: $\kt < 0$ gives $\Lambda_{\mathrm{ent}} > 0$ (positive cosmological-constant-like
term, repulsive curvature).

\textbf{Modified Einstein equations.}
\begin{equation}
  G_{\mu\nu}
  \;\equiv\;
  R_{\mu\nu} - \tfrac{1}{2}R\,g_{\mu\nu}
  \;=\;
  \frac{8\pi G}{c^4}\,T^{\mathrm{total}}_{\mu\nu}
  \;=\;
  \frac{8\pi G}{c^4}\bigl(T^{\mathrm{matter}}_{\mu\nu} - \Lambda_{\mathrm{ent}}\,g_{\mu\nu}\bigr).
  \label{eq:EFE}
\end{equation}

\textbf{Contracted Bianchi identity.}  The second contracted Bianchi identity
$\nabla^\mu G_{\mu\nu} \equiv 0$ forces:
\begin{equation}
  \nabla^\mu T^{\mathrm{total}}_{\mu\nu} = 0
  \;\;\Longrightarrow\;\;
  \nabla^\mu T^{\mathrm{matter}}_{\mu\nu}
  = -\nabla^\mu T^{\mathrm{ent}}_{\mu\nu}
  = \nabla_\nu\Lambda_{\mathrm{ent}}
  = \kt\,\frac{c^4}{8\pi G\kB\ln 2}\,\nabla_\nu\Sent.
  \label{eq:bianchi}
\end{equation}
This is an \emph{exact} consequence of the field equations~\eqref{eq:EFE}.  It is not an
additional assumption.  It forces the interpretation adopted in the white paper:
the right-hand side of Eq.~\eqref{eq:bianchi} is the energy-momentum current flowing
from the entanglement sector into the matter sector (or vice versa, depending on the sign
of $\nabla_\nu\Sent$).

\begin{expertflag}
Equation~\eqref{eq:bianchi} implies a physical prediction that the white paper
has not yet confronted explicitly.  If matter is a point mass $m$ at rest, its equation
of motion in this background is
\begin{equation}
  m\,\frac{Du^\mu}{d\tau}
  = \kt\,\frac{c^4}{8\pi G\kB\ln 2}\,\frac{\nabla^\mu\Sent}{T^{\mathrm{matter}}_{00}/c^2},
  \label{eq:eom}
\end{equation}
where $D/d\tau$ is the covariant proper-time derivative.  This is a classical fifth force
proportional to $\nabla\Sent$ and independent of composition.  The atom interferometry
bound $\eta < 10^{-12}$~\cite{Asenbaum2020} bounds \emph{composition-dependent} forces,
but if two test masses are placed in regions of different $\nabla\Sent$,
Eq.~\eqref{eq:eom} predicts different accelerations even without a new force law.
This observable is not ruled out by current experiments and should be
stated explicitly in the paper.
\end{expertflag}

\subsection{Holographic regulator steps: every assumption labeled}

The estimate $\kt = -1/4$ requires two assumptions beyond Jacobson's theorem.
Both are labeled here with \textbf{[HRH]} (Holographic Regulator Hypothesis):

\begin{enumerate}[label=\textbf{Step \arabic*}]
  \item \textbf{[Jacobson, established]} Apply $\delta Q = T_U\,dS_{BH}$ to a local
    Rindler horizon with $T_U = \hbar a/(2\pi c\kB)$ and
    $dS_{BH} = (\kB c^3/4G\hbar)\,dA$.  This gives $G_{\mu\nu} = 8\pi G c^{-4}T^{\mathrm{matter}}_{\mu\nu}$.

  \item \textbf{[HRH-1]} For a volume with entanglement entropy density $\Sent$ (bit\,m$^{-3}$),
    propose an additional entropy contribution to the horizon element $dA$:
    \begin{equation}
      dS_{\mathrm{ent}} = \frac{\Sent}{\kB}\cdot\frac{\lP\,dA}{4\lP}
      = \frac{\Sent}{4\kB}\,dA.
      \label{eq:dSent}
    \end{equation}
    \textit{Assumption:} the volume element behind the horizon is exactly one Planck length
    deep ($dV = \lP\,dA$), so the full volume entropy density collapses to an area density
    via $\Sent\,dV \to \Sent\,\lP\,dA$.  This regulation has no rigorous justification
    for non-horizon laboratory volumes.  It is valid at causal horizons (where
    Bekenstein's bound enforces $S \leq A/4\lP^2$) but not for bulk systems with volume-law
    entanglement.

  \item \textbf{[HRH-2]} Add $T_U\,dS_{\mathrm{ent}}$ to the Clausius heat flux.
    Following Jacobson's identification of heat with stress-energy flux gives:
    \begin{equation}
      T^{\mathrm{ent}}_{\mu\nu} = -\frac{c^4}{32\pi G}\Sent\,g_{\mu\nu},
      \label{eq:Tent-derived}
    \end{equation}
    which, matched to Eq.~\eqref{eq:Tent}, yields $\kt = -1/4$.

    \textit{Assumption (implicit):} The Clausius identification $\delta Q \leftrightarrow
    T^{\mathrm{eff}}_{\mu\nu}k^\mu d\Sigma^\nu$ extends to the entanglement contribution.
    This requires that the entanglement entropy gradient contributes thermodynamically
    to the local horizon in the same way as matter energy-momentum.  No proof exists.
\end{enumerate}

The result $\kt = -1/4$ is therefore contingent on both HRH-1 and HRH-2.  Removing either
assumption leaves the coupling as a free parameter.

\subsection{Minimal addition to the source term for a natural conservation story}

The current source term $T^{\mathrm{ent}}_{\mu\nu} = -\Lambda_{\mathrm{ent}}(x)g_{\mu\nu}$
is the simplest possible form (cosmological-constant structure with a varying prefactor).
Its conservation equation~\eqref{eq:bianchi} forces energy-momentum exchange with matter,
which is physically plausible but lacks an action principle.

\textbf{Proposed minimal extension: kinetic scalar-field term.}

Promote $\Sent$ to a dynamical scalar field $\phi$ and write:
\begin{equation}
  T^{\mathrm{ent}}_{\mu\nu}
  \;=\;
  \alpha_\phi\!\left(\nabla_\mu\phi\,\nabla_\nu\phi
  - \tfrac{1}{2}g_{\mu\nu}(\nabla\phi)^2\right)
  - \Lambda_\phi(\phi)\,g_{\mu\nu},
  \label{eq:Tent-full}
\end{equation}
where $\alpha_\phi$ is a positive constant and $\Lambda_\phi(\phi) = -\kt(c^4/8\pi G\kB\ln 2)\phi$.
This has a natural action:
\begin{equation}
  S_{\mathrm{ent}}[\phi,g]
  = \int d^4x\,\sqrt{-g}
    \left[
      -\frac{\alpha_\phi}{2}(\nabla^\mu\phi\,\nabla_\mu\phi)
      + \Lambda_\phi(\phi)
    \right].
  \label{eq:action}
\end{equation}
Varying with respect to $\phi$ gives the scalar field equation of motion:
\begin{equation}
  \alpha_\phi\,\Box\phi + \frac{d\Lambda_\phi}{d\phi}
  = \alpha_\phi\,\Box\phi + \kt\frac{c^4}{8\pi G\kB\ln 2}
  = 0,
  \label{eq:phi-eom}
\end{equation}
which relates the propagation of $\phi$ to the entanglement source density.
Varying with respect to $g^{\mu\nu}$ gives Eq.~\eqref{eq:Tent-full} as the entanglement
stress-energy.  The \emph{combined} system $(g_{\mu\nu}, \phi)$ satisfies
$\nabla^\mu T^{\mathrm{total}}_{\mu\nu} = 0$ automatically from the action principle.

\textbf{Physical interpretation.}  In the limit $\alpha_\phi \to 0$, $\phi = \Sent$ is fixed
by external conditions (the quantum state of the system), and Eq.~\eqref{eq:Tent-full}
reduces to the white paper's current ansatz.  For $\alpha_\phi > 0$, $\phi$ propagates as a
massless scalar and can exchange energy-momentum with matter via the coupling in
Eq.~\eqref{eq:bianchi} without violating any conservation law.

\begin{remark}
The scalar $\phi$ is not a new fundamental field.  It is an effective field representing the
coarse-grained entanglement entropy density of the quantum system --- an emergent degree of
freedom, in the same spirit that thermodynamic quantities are emergent from microphysics.
The action~\eqref{eq:action} should be understood as an effective field theory valid below
the Planck scale, not a UV-complete theory.  Whether $\phi$ can be derived from a fundamental
quantum gravity theory (e.g., via holographic renormalization group) is Open Problem~10.
\end{remark}

\textbf{Conservation in the extended theory.}  With the action~\eqref{eq:action}, the
stress-energy of the full system satisfies:
\begin{align}
  \nabla^\mu T^{\mathrm{matter}}_{\mu\nu}
  &= -\nabla^\mu T^{\mathrm{ent}}_{\mu\nu}, \label{eq:cons1} \\
  \nabla^\mu T^{\mathrm{ent}}_{\mu\nu}
  &= -\alpha_\phi\,(\Box\phi)\,\nabla_\nu\phi
    + \nabla_\nu\Lambda_\phi
    = \alpha_\phi(\alpha_\phi^{-1}\nabla_\nu\Lambda_\phi)\cdot(-\Box\phi)
    + \nabla_\nu\Lambda_\phi, \label{eq:cons2}
\end{align}
which vanishes on-shell (when $\phi$ satisfies its equation of motion~\eqref{eq:phi-eom}).
This is the correct, covariant conservation story.

% ─────────────────────────────────────────────────────────────────────────────
\section{OP3: A Concrete Lindblad Model for Screening}
\label{sec:op3}

\subsection{The precise definition of $S_{\mathrm{ent}}$ that must be used}

Before writing any Lindblad model, the white paper must specify exactly which entropy
enters $T^{\mathrm{ent}}_{\mu\nu}$.  There are three candidates and they have
\emph{opposite} behaviors under decoherence:

\begin{enumerate}
  \item \textbf{Total von Neumann entropy of the quantum system:}
    $S_{\mathrm{total}} = -\kB\tr(\rho\ln\rho)$.
    Under decoherence, this \emph{increases} (purity is lost).
    Using this would make decoherence \emph{increase} the gravitational source --- physically
    unnatural and inconsistent with the intuition that entanglement sources the effect.

  \item \textbf{System-environment entanglement entropy:}
    $S_{\mathrm{SE}} = -\kB\tr(\rho_S\ln\rho_S)$ where $\rho_S = \tr_E(\rho_{SE})$.
    This also increases under decoherence.  Same problem.

  \item \textbf{Internal (inter-qubit) entanglement entropy:}
    The von Neumann entropy of the reduced state of a subsystem traced over the other
    subsystems of the \emph{same} quantum device, treating the environment as external.
    Concretely, for a bipartition $A|B$ of the device:
    $S_{\mathrm{int}} = -\kB\tr(\rho_A\ln\rho_A)$ where $\rho_A = \tr_B\rho_{\mathrm{device}}$.
    Under decoherence, $\rho_{\mathrm{device}} \to$ a nearly classically correlated state,
    and $S_{\mathrm{int}} \to 0$ (the system becomes product after tracing).
    This \emph{decreases} under decoherence and is what the white paper implicitly intends.
\end{enumerate}

\begin{expertflag}
The white paper has not stated which of the three candidates it uses.  Candidate~(3) is
the physically correct choice --- it is the quantity that (a) decreases under decoherence,
(b) is zero for a fully decohered state, and (c) is what experimentalists can
verify via entanglement witnesses.  All subsequent analysis in OP3 uses $\Sent = S_{\mathrm{int}}$:
the bipartite (or multi-partite) entanglement entropy of the internal degrees of freedom of
the quantum device, tracing over the environment.  This definition must be added explicitly
to the white paper's Section~2.
\end{expertflag}

\subsection{Lindblad model for 50-qubit superconducting transmon qubits}

\textbf{Platform choice.}  We model a chain of $N = 50$ fixed-frequency transmon qubits,
which is the architecture currently operated by IBM Quantum and Google.  This choice is made
for concreteness; the hierarchy developed below applies qualitatively to all NISQ platforms.

\textbf{System Hamiltonian.}  In the rotating frame, in the two-level (qubit) approximation:
\begin{equation}
  H_{\mathrm{sys}}
  = \sum_{k=1}^{N}\frac{\omega_k}{2}\sigma^z_k
  + \sum_{\langle j,k\rangle}J_{jk}\left(\sigma^+_j\sigma^-_k + \mathrm{h.c.}\right),
  \label{eq:H-transmon}
\end{equation}
where $\omega_k/2\pi \in [4.5, 5.0]$\,GHz are qubit frequencies and
$J_{jk}/2\pi \sim 2$--$10$\,MHz are nearest-neighbor exchange couplings mediated by
coplanar microwave resonators.

\textbf{Lindblad operators and rates.}  The dominant environmental channels are:

\begin{enumerate}[label=(\roman*)]
  \item \textit{Energy relaxation} (T1 process):
    $L_{1,k} = \sqrt{\gamma_1}\,\sigma^-_k$, rate $\gamma_1 = 1/T_1$.
    State of the art: $T_1 \sim 100$--$500\,\mu$s (Google/IBM, 2025).

  \item \textit{Pure dephasing} (T2 process):
    $L_{\phi,k} = \sqrt{\gamma_\phi/2}\,\sigma^z_k$, rate $\gamma_\phi = 1/T_\phi$
    where $T_\phi^{-1} = T_2^{-1} - T_1^{-1}/2$.
    State of the art: $T_2 \sim 100\,\mu$s, giving $T_\phi \sim 200\,\mu$s.

  \item \textit{Thermal excitation} (small at $T_{\mathrm{fridge}} \sim 15$\,mK):
    $L_{\mathrm{exc},k} = \sqrt{\bar{n}\gamma_1}\,\sigma^+_k$,
    where $\bar{n} = (e^{\hbar\omega_k/\kB T_{\mathrm{fridge}}}-1)^{-1} \lesssim 10^{-4}$.
    Negligible.
\end{enumerate}

The master equation is:
\begin{equation}
  \dot{\rho}
  = -\frac{i}{\hbar}[H_{\mathrm{sys}},\rho]
  + \sum_{k=1}^{N}\sum_{\alpha\in\{1,\phi\}}\gamma_\alpha\!
    \left(L_{\alpha,k}\,\rho\,L_{\alpha,k}^\dagger
    - \tfrac{1}{2}\{L_{\alpha,k}^\dagger L_{\alpha,k},\rho\}\right).
  \label{eq:lindblad}
\end{equation}

\subsection{Which correlators couple to the gravitational source}

The gravitational source $\Sent = S_{\mathrm{int}}$ depends on the full reduced density
matrix $\rho_S(t) = \tr_E\rho_{SE}(t)$.  For computational purposes, the most tractable
proxy is the \emph{single-site entropy}:
\begin{equation}
  s_k(t) = -\kB\tr_k\!\left[\rho_k(t)\ln\rho_k(t)\right],
  \quad
  \rho_k = \tr_{\neq k}\rho_S,
  \label{eq:single-site}
\end{equation}
where $\tr_{\neq k}$ traces over all qubits except qubit $k$.  For a pure state,
$s_k$ is the entanglement entropy between qubit $k$ and the rest.

For a GHZ state $|\text{GHZ}\rangle = (|0\rangle^{\otimes N}+|1\rangle^{\otimes N})/\sqrt{2}$,
$s_k = \kB\ln 2$ for every $k$.  Under qubit $k$'s dephasing at rate $\gamma_\phi$:
\begin{equation}
  s_k(t) = \kB\,h\!\left(\frac{1 + e^{-\gamma_\phi t}}{2}\right),
  \label{eq:sk-decay}
\end{equation}
where $h(p) = -p\ln p - (1-p)\ln(1-p)$ is binary entropy.  However, a single-qubit
dephasing event \emph{also} collapses the GHZ entanglement between all other qubits:
the GHZ correlation $|\langle 0^{\otimes N}|\rho|1^{\otimes N}\rangle|$ decays as
$e^{-\gamma_\phi t}$ under a \emph{single-qubit} dephasing with rate $\gamma_\phi$,
but under \emph{any-qubit} dephasing the full GHZ coherence decays at rate $N\gamma_\phi$.

Therefore, the integrated entanglement entropy for GHZ during a measurement protocol of
duration $\tau$:
\begin{equation}
  \Ment^{\mathrm{GHZ}}
  \approx N\kB\ln 2 \int_0^\tau e^{-N\gamma_\phi t}\,dt
  = \frac{\kB\ln 2}{\gamma_\phi}\left(1 - e^{-N\gamma_\phi\tau}\right).
  \label{eq:Ment-GHZ}
\end{equation}
Maximized at $\tau \to \infty$: ${\Ment}_{\mathrm{GHZ,max}} = \kB\ln 2/\gamma_\phi$,
independent of $N$ (the GHZ state's total entanglement is lost on the timescale
$1/(N\gamma_\phi)$, which decreases with $N$).

For a \textbf{cluster state}, which has short-range entanglement (each qubit entangled with
its neighbors), the situation is more favorable.  Under single-qubit dephasing:
\begin{equation}
  {\Ment}_{\mathrm{cluster}}
  \approx N\kB\ln 2 \int_0^\tau e^{-\gamma_\phi t}\,dt
  = \frac{N\kB\ln 2}{\gamma_\phi}\left(1 - e^{-\gamma_\phi\tau}\right).
  \label{eq:Ment-cluster}
\end{equation}
The cluster state's total entanglement grows with $N$ because its local entanglement is
protected from global collapse.  This suggests cluster states are superior to GHZ states
for maximizing $\Ment$ under realistic decoherence.

\textbf{Key finding:} The choice of entangled state significantly affects the effective
$\Ment$, and thus $\asc$:
\begin{equation}
  \frac{{\Ment}_{\mathrm{cluster}}}{\Ment^{\mathrm{GHZ}}}
  \approx N\left(1 - e^{-\gamma_\phi\tau}\right)
  \gg 1 \quad \text{for } N \gg 1.
  \label{eq:ratio}
\end{equation}
The experiment should use cluster states (or matrix product states with finite correlation
length), not GHZ states, for maximum signal.

\subsection{Hierarchy of regimes: qualitative expectations for $\asc$}

We define the effective screening factor as the ratio of actual to ideal integrated
entanglement entropy:
\begin{equation}
  \asc \equiv \frac{\Ment^{\mathrm{actual}}}{\Ment^{\mathrm{ideal}}}.
  \label{eq:asc-def}
\end{equation}

\textbf{Regime I: Few-qubit ($N \lesssim 20$), short circuits ($d \lesssim 10$ gate layers).}
\begin{itemize}
  \item Gate fidelity: $F_{\mathrm{2Q}} \geq 99.5\%$; circuit error per qubit $\varepsilon \lesssim 0.02$
  \item Coherence: $T_2 \gg N\tau_{\mathrm{gate}}$; negligible decoherence during circuit
  \item Entanglement depth: achievable up to $k = N$ (full entanglement depth)
  \item $\asc$ estimate: $0.1 \lesssim \asc \lesssim 0.5$ (significant but not catastrophic suppression)
  \item $\Ment$: small (few qubits), signal correspondingly small
\end{itemize}

\textbf{Regime II: Mesoscopic ($N = 50$--$100$), moderate circuits ($d \sim 20$--$50$ layers).}
\begin{itemize}
  \item Cluster state preparation requires $d \sim N/2 \sim 25$--$50$ CZ gate layers
  \item Decoherence per layer: $\varepsilon \sim N\gamma_\phi\tau_{\mathrm{gate}} \sim 50\times10^4\times10^{-7} \approx 5\times10^{-3}$ per qubit per layer
  \item Total entanglement decay: cluster state entanglement decays at $\sim d\varepsilon \sim 0.1$--$0.25$ (10--25\% loss per circuit)
  \item $\asc$ estimate (cluster state): $e^{-d\varepsilon} \sim 10^{-1}$ to $10^{-2}$
  \item $\asc$ estimate (GHZ state): $e^{-N\gamma_\phi\tau} \sim 10^{-3}$ to $10^{-5}$
  \item This is the \textbf{optimal regime}: large enough $\Ment$, small enough suppression
\end{itemize}

\textbf{Regime III: Macroscopic ($N \gtrsim 10^3$), deep circuits or long coherence time.}
\begin{itemize}
  \item GHZ-type global entanglement: $\asc \sim e^{-N\gamma_\phi\tau} \to 0$ rapidly
  \item Local entanglement (cluster, MPS): $\asc$ may be $\mathcal{O}(0.01)$ but
    $\Ment/N^{\alpha}$ saturates as correlation length limits
  \item Macroscopic BEC: $R/\xi \sim 10$--$100$, Gaussian screening gives
    $e^{-(R/\xi)^2} \sim e^{-100}$ to $e^{-10000}$: effectively zero
  \item Verdict: not a viable near-term target for NISQ-based entanglement-gravity search
\end{itemize}

\begin{openprobbox}
\textbf{OP3 Lindblad milestone path.}
\begin{enumerate}
  \item \textit{Month 1--2}: For a 50-qubit transmon chain with Hamiltonian~\eqref{eq:H-transmon}
    and Lindblad operators above, numerically solve~\eqref{eq:lindblad} for $N = 5$, $10$, $20$
    qubits with cluster state initial conditions.  Compute $s_k(t)$ and $\Ment(t)$ as functions
    of $\gamma_\phi\tau$.  This produces a calibrated $\asc(\gamma_\phi, N, \tau)$ function for
    the cluster state.
  \item \textit{Month 2--4}: Extend the Lindblad model to include gate errors (coherent and
    incoherent) using gate set tomography~\cite{Blume2017} data from actual hardware.  Replace
    ideal unitary gates with their Kraus-map representations.  Recompute $\Ment$ including
    gate-level decoherence.
  \item \textit{Month 4--6}: Scale the numerical solution to $N = 50$ using tensor network
    (MPO/TDVP) methods.  Identify the circuit depth $d^*$ that maximizes
    $\Ment^{\mathrm{total}} = \int_0^\tau \sum_k s_k(t)\,dt$.
  \item \textit{Month 6--12}: Analytically bound $\asc$ using the quantum data processing
    inequality and the monotonicity of relative entropy under CPTP maps.  Establish whether
    $\asc$ is bounded away from zero for any circuit on any platform, or whether it
    necessarily vanishes in the thermodynamic limit.
\end{enumerate}
The outcome of step 4 is the most important: if $\asc \to 0$ for all $N \to \infty$,
the framework's experimental program must be confined to the few-qubit regime.
\end{openprobbox}

% ─────────────────────────────────────────────────────────────────────────────
\section{NISQ Experiment: Publishable Proposal}
\label{sec:experiment}

\subsection{Platform: 2D reconfigurable neutral-atom array}

\textbf{Architecture.}
A $6 \times 6 = 36$ qubit array of $^{87}$Rb atoms held in optical tweezer traps
at lattice spacing $a = 5\,\mu$m.  Two-qubit gates via Rydberg blockade using
$|r\rangle = |n\!=\!70, S_{1/2}\rangle$ (blockade radius $r_b \approx 8\,\mu$m for
nearest-neighbor coupling, $r_b/a \approx 1.6$, ensuring blockade on neighbors only).

\textbf{Hilbert space.}
Qubit basis: $|0\rangle = |5S_{1/2}, F=1, m_F=0\rangle$,
$|1\rangle = |5S_{1/2}, F=2, m_F=0\rangle$ (clock transition at $\omega_{hf}/2\pi = 6.835$\,GHz).
Full Hilbert space: $\mathcal{H} = (\mathbb{C}^2)^{\otimes 36}$.

\textbf{Hamiltonian} (rotating frame, Rydberg dressing on for entangling gates):
\begin{equation}
  H
  = \sum_{k=1}^{36}\!\left(
      \frac{\Omega_k}{2}\sigma^x_k - \frac{\Delta_k}{2}\sigma^z_k
    \right)
  + \sum_{k < l}\!\frac{C_6}{|\mathbf{r}_k - \mathbf{r}_l|^6}\,
    \hat{n}^{(r)}_k\,\hat{n}^{(r)}_l,
  \label{eq:H-rydberg}
\end{equation}
where $C_6/h = 862\,\text{GHz}\cdot\mu\text{m}^6$ for the $|70S_{1/2}\rangle$ state,
$\Omega_k$ and $\Delta_k$ are Rabi frequency and detuning from a global or site-addressed
laser, and $\hat{n}^{(r)}_k = |r\rangle\langle r|_k$ is the Rydberg population operator.
During gates, $\Omega_k/2\pi \approx 3$\,MHz and $\Delta_k = 0$.

\textbf{Decoherence parameters.}
\begin{itemize}
  \item Rydberg lifetime: $\tau_r \approx 0.3$\,ms at $T_{\mathrm{atom}} \sim 5\,\mu$K
  \item CZ gate time: $\tau_{\mathrm{CZ}} \approx 200$\,ns (limited by $1/\Omega$)
  \item Single-qubit gate time: $\tau_{\mathrm{1Q}} \approx 500$\,ns
  \item Coherence time (clock states): $T_2 > 1$\,s (magnetic-field-insensitive transition)
  \item Readout fidelity: $F_{\mathrm{read}} \approx 99.5\%$ (site-resolved fluorescence)
\end{itemize}

\subsection{Branch A: High-entanglement states}

\textbf{State A1 (2D cluster state).}  Prepare via:
\begin{enumerate}
  \item Initialize all qubits in $|+\rangle^{\otimes 36}$
  \item Apply CZ$_{kl}$ to all nearest-neighbor pairs (two layers: horizontal + vertical)
  \item Result: 2D cluster state with entanglement depth $k = 36$ and
    bipartite entanglement $S_{A|B} \approx \min(|A|, |B|)\ln 2$ for any cut $A|B$
\end{enumerate}
Circuit depth: 2 layers of CZ gates + 1 layer of Hadamards.
Total time: $\tau_{\mathrm{prep}} \approx 2 \times 200\,\text{ns} + 500\,\text{ns} \approx 1\,\mu$s.
Entanglement entropy: $\Ment^{\mathrm{tot}} \approx 36\kB\ln 2 \approx 25\,\kB$.

\textbf{State A2 (Checkerboard GHZ).}  Divide the 6$\times$6 array into two sublattices
(checkerboard pattern); prepare an 18-qubit GHZ state on each sublattice.
Less fragile than full-array GHZ; $\Ment^{\mathrm{tot}} \approx 2\kB\ln 2$.

\subsection{Branch B: Low-entanglement reference states}

\textbf{State B1 (Product state, identical energy).}
Prepare $|+\rangle^{\otimes 36}$ with the same laser schedule as Branch~A but with all
CZ gates replaced by identity operations (zero-detuning pulses that return to the
computational basis without entangling).  Energy difference from A: at most
$\sim N\langle J_{jk}\rangle \sim 36 \times 5$\,MHz $\sim 10$\,kHz --- must be characterized.

\textbf{State B2 (Two-qubit entangled product).}
Prepare Bell pairs $|\Phi^+\rangle^{\otimes 18}$ (all 18 nearest-neighbor pairs
entangled but no multipartite entanglement).  $\Ment^{\mathrm{tot}} \approx 18\kB\ln 2$
(different from A1 by a factor $\sim 2$, offering an intermediate comparison point).

\subsection{Observable: Ramsey phase differential}

\textbf{Protocol.}
\begin{enumerate}
  \item Prepare Branch A (or B) state
  \item Apply a global $\pi/2$ rotation around $\hat{x}$: maps $|0\rangle \to |+\rangle$,
    $|1\rangle \to |-\rangle$ on the clock transition
  \item Free precession for time $T_{\mathrm{Ramsey}}$
  \item Apply a second $\pi/2$ rotation; measure population in $|1\rangle$
  \item The accumulated phase $\phi = \omega_{\mathrm{eff}}\,T_{\mathrm{Ramsey}}$ encodes
    any shift in the effective hyperfine frequency
\end{enumerate}

\textbf{Signal structure.}  A coupling of the form $T^{\mathrm{ent}}_{\mu\nu} = -\Lambda_{\mathrm{ent}}g_{\mu\nu}$
modifies the local metric as $\delta g_{00} \approx -2\Phi_{\mathrm{ent}}/c^2$ where
$\Phi_{\mathrm{ent}}$ is the Newtonian potential of the entanglement source.  This produces
a gravitational redshift of the hyperfine transition:
\begin{equation}
  \frac{\delta\omega_{\mathrm{hf}}}{\omega_{\mathrm{hf}}}
  = \frac{\Phi_{\mathrm{ent}}}{c^2}
  = -\frac{\kteff c^2}{8\pi G\kB\ln 2}\frac{\Ment}{R_{\mathrm{eff}}^2},
  \label{eq:frequency-shift}
\end{equation}
where $R_{\mathrm{eff}}$ is the effective distance scale set by the qubit array geometry.

\textbf{What this actually measures.}
The differential phase between Branch~A and Branch~B accumulated over $T_{\mathrm{Ramsey}}$:
\begin{equation}
  \Delta\phi_{AB}
  = (\omega^A_{\mathrm{eff}} - \omega^B_{\mathrm{eff}})\,T_{\mathrm{Ramsey}}
  = \frac{\delta\omega^{AB}_{\mathrm{hf}}}{\omega_{\mathrm{hf}}}\,\omega_{\mathrm{hf}}\,T_{\mathrm{Ramsey}}.
  \label{eq:phase-diff}
\end{equation}
This quantity bounds the combination
\begin{equation}
  \left|\kteff\right| \cdot \frac{\Ment^A - \Ment^B}{R_{\mathrm{eff}}^2}
  < \frac{|\Delta\phi_{AB}|_{\mathrm{sens}}}{T_{\mathrm{Ramsey}}}
    \cdot \frac{8\pi G\kB\ln 2}{c^2\,\omega_{\mathrm{hf}}},
  \label{eq:bound}
\end{equation}
without requiring knowledge of $\asc$ or the Planck-scale renormalization.
The right-hand side is a pure number set by the experimental sensitivity.

For current state-of-the-art Ramsey spectroscopy on Rb clock transitions:
sensitivity $\delta(\delta\omega)/\omega \sim 10^{-14}$--$10^{-16}$ per shot in optimal
conditions (limited by quantum projection noise at $\delta\phi \sim 1/\sqrt{N_{\mathrm{atoms}}}$
per shot, and by measurement time).  After $10^4$ shots: $\delta(\Delta\phi) \sim 10^{-3}$\,rad.

\subsection{Systematic error budget and fatal vs.\ non-fatal errors}

\begin{table}[h]
\centering\small
\caption{Systematic error budget for the neutral-atom Ramsey protocol.}
\label{tab:systematics}
\begin{tabular}{@{}p{4.5cm}p{3.5cm}p{3.5cm}p{2cm}@{}}
\toprule
Systematic & Physical origin & Mitigation & Fatal? \\
\midrule
Differential mean-field energy shift
& $\sum_{kl}J_{kl}(\langle\hat{n}_k\hat{n}_l\rangle_A - \langle\hat{n}_k\hat{n}_l\rangle_B)$
& Match two-body correlators between A and B (e.g., use B2 with same $\langle\hat{n}\hat{n}\rangle$)
& \textbf{Yes} if uncontrolled \\
\addlinespace
Differential AC Stark shift from gate pulses
& Gate laser amplitude $\propto\Omega^2/\Delta_{\mathrm{laser}}$
& Zero-area pulse sequences; same laser schedule for A and B
& No \\
\addlinespace
Magnetic field fluctuations
& $\partial B/\partial t \neq 0$
& Clock transition insensitive to first order; active shielding
& No \\
\addlinespace
Differential trap displacement (A vs B)
& Different motional state for entangled atoms
& Common-mode rejection from differential measurement; same trapping geometry
& No \\
\addlinespace
SPAM errors in entanglement verification
& Readout fidelity $< 1$
& Entanglement witnesses with SPAM calibration (GST)
& No \\
\addlinespace
Atom loss during circuit
& Rydberg decay, collision loss
& Monitor atom number; postselect on full count
& No (postselect) \\
\addlinespace
Differential motional heating
& Heating rate $\dot{n}_{\mathrm{phonon}}$
& Sideband cooling; matched motional temperature for A and B
& No \\
\bottomrule
\end{tabular}
\end{table}

\textbf{The fatal systematic in detail.}
The entangled state (A1, 2D cluster) has different two-body correlators
$\langle\hat{n}_j\hat{n}_k\rangle$ than the product state (B1).  For the cluster state,
$\langle\hat{n}_j\hat{n}_k\rangle_{\mathrm{cluster}} = 1/4$ for all pairs
(since each qubit is in a $|+\rangle$ state locally).  For $|+\rangle^{\otimes N}$,
$\langle\hat{n}_j\hat{n}_k\rangle_{\mathrm{product}} = 1/4$ as well.  So this specific
pair of states has \emph{identical} two-body correlators.  This is why cluster states
are preferred over GHZ for this comparison: the differential mean-field energy is zero.

However, for GHZ vs.\ product: $\langle\hat{n}_j\hat{n}_k\rangle_{\mathrm{GHZ}} = 1/4$
but the GHZ has non-trivial $N$-body correlators that create different collective
energy shifts.  This must be verified numerically for the specific Hamiltonian.

\textbf{Null-result interpretation.}
If $\Delta\phi_{AB} = 0$ at the level of experimental sensitivity, this rules out:
\begin{equation}
  |\kteff|\cdot\frac{\Delta\Ment}{R_{\mathrm{eff}}^2}
  > \frac{|\Delta\phi|_{\mathrm{sens}}\,8\pi G\kB\ln 2}{c^2\,\omega_{\mathrm{hf}}\,T_{\mathrm{Ramsey}}}.
\end{equation}
With $R_{\mathrm{eff}} \sim 3a = 15\,\mu$m, $\Delta\Ment \sim 18\kB\ln 2$,
$T_{\mathrm{Ramsey}} = 0.1$\,s, $|\Delta\phi|_{\mathrm{sens}} \sim 10^{-3}$\,rad:

$|\kteff| < \frac{10^{-3} \times 8\pi \times 6.67\times10^{-11} \times 9.57\times10^{-24}}
  {(3\times10^8)^2 \times 6.84\times10^9 \times 0.1}
  \times \frac{(1.5\times10^{-5})^2}{18 \times 9.57\times10^{-24}}
\approx 10^{-47}$.

This is an extraordinarily small bound on $\kteff$, confirming OP10: the coupling prefactor
$c^4/(8\pi G\kB\ln 2)$ is so large that even with $\asc \sim 10^{-2}$, the experiment
constrains $\kt$ to be far below $10^{-15}$ --- meaning the falsification criterion F1--F5
is much harder to reach than the white paper suggests at face value.

\begin{expertflag}
The experimental bound on $|\kteff|$ from the proposed NISQ experiment is
approximately $10^{-47}$, not the $10^{-15}$ stated in the falsification criterion.
This reflects the enormous ratio $c^4/(8\pi G \kB\ln 2) \approx 5\times10^{65}$ in SI units.
The gap between $10^{-47}$ and $10^{-15}$ (a factor of $10^{32}$) is precisely OP10:
if the Planck-scale renormalization suppresses $\kt$ to $\sim 10^{-47}$ at laboratory scales,
the framework is self-consistent and both bounds agree.  If $\kt$ is renormalized to a value
larger than $10^{-47}$ (e.g., $10^{-15}$), the proposed experiment has already falsified it.
This should be stated explicitly in the paper's experimental section.
\end{expertflag}

% ─────────────────────────────────────────────────────────────────────────────
\section{Mock Referee Report and Three-Paper Sequence}
\label{sec:referee}

\subsection{Mock referee report}

\begin{mdframed}[backgroundcolor=boxgray,linewidth=1pt]
\textbf{Referee Report}\\
\textit{Journal:} Physical Review D (or Classical and Quantum Gravity)\\
\textit{Manuscript:} Entanglement Entropy Density as a Gravitational Source

\medskip\noindent\textbf{SUMMARY}

The manuscript presents a phenomenological extension of Jacobson's thermodynamic gravity
in which entanglement entropy density acts as an additional gravitational source term.
The dimensional consistency of the modification is established, and the paper is
admirably honest about where derivations end and hypotheses begin.  The resolution of the
radial scaling question (OP4) and the treatment of the Bianchi identity are genuine
improvements over prior drafts.

\medskip\noindent\textbf{MAJOR STRENGTHS}

(1) \textit{Epistemic clarity.}  The separation between established results (dimensional
consistency, radial scaling), motivated hypotheses (Planck regulator, $\kt = -1/4$), and
unresolved parameters ($\asc$, renormalization) is one of the clearest in speculative
gravity phenomenology literature.  This alone makes the paper worth reading.

(2) \textit{OP4 resolution.}  The derivation $\Delta a(R) \propto \Ment(<\!R)/R^2$
from the modified Poisson equation is correct and removes a structural inconsistency
from prior work.

(3) \textit{Falsification orientation.}  The F1--F5 criterion is quantitative and
platform-independent.  The paper correctly notes that null results are scientifically
valuable even before $\asc$ is resolved.

\medskip\noindent\textbf{MAJOR CONCERNS (revisions required)}

(C1) \textit{Definition of $\Sent$ is missing.}
The manuscript uses $\Sent$ throughout but never states precisely which entropy it is.
Three options exist (total von Neumann, system-environment, internal inter-qubit) with
opposite behaviors under decoherence.  The analysis requires candidate~(3) (internal
entanglement entropy that decreases under decoherence); this must be stated explicitly.
Without it, the sign of $\asc$'s effect on the gravitational source is ambiguous.

(C2) \textit{The fifth-force prediction from the Bianchi identity is not confronted.}
Equation~\eqref{eq:bianchi} predicts a classical force on matter proportional to $\nabla\Sent$.
For a single qubit near an entangled register, this force is not composition-dependent
(unlike MICROSCOPE's constraint) but is entanglement-density-dependent.  The paper must
either (a) bound this force from existing precision measurements, or (b) state explicitly
that it is an unconstrained prediction.  Claiming MICROSCOPE does not constrain $\kt$ is
correct but incomplete: MICROSCOPE tests composition-dependence, not entanglement-dependence.

(C3) \textit{The signal magnitude is inconsistent with the falsification criterion.}
The falsification criterion (F2: $\Delta p < 10^{-6}$\,Pa) and the proposed NISQ
experiment produce a bound on $|\kteff|$ that is approximately $10^{-47}$, not $10^{-15}$.
Either the falsification criterion should be tightened to match realistic experiments,
or the criterion should be rephrased as bounding $|\kteff|\cdot\Ment/R^2$ directly
without implying anything about $\kt$ alone.

(C4) \textit{OP3 is not specific enough to constitute an open problem statement.}
``Derive $\asc$ from a Lindblad master equation'' is a research direction, not an open
problem.  A publishable open problem statement must specify: the platform, the specific
correlators that couple to $\Sent$, the Lindblad operators, and what a ``solution''
would consist of.  The milestone path outlined in Section~2.4 of this analysis
is the kind of specificity required.

(C5) \textit{The conservation interpretation (entanglement-sector exchange) is adopted
without justification.}  Of the three interpretations (fifth force, sector exchange,
incomplete source term), interpretation~(2) is chosen as ``more natural'' without
argument.  A natural justification requires an action principle, which the manuscript
lacks.  The paper should either (a) present the scalar-field effective action
(Eq.~\eqref{eq:action}) as minimal supporting structure, or (b) explicitly state that
the conservation story is an ansatz awaiting an action derivation.

\medskip\noindent\textbf{SPECIFIC REQUIRED REVISIONS}

(R1) Add explicit definition: $\Sent = S_{\mathrm{int}}(x)$ = bipartite entanglement entropy
density of the internal quantum degrees of freedom, tracing over the environment.  Clarify
that this decreases under decoherence.

(R2) Add a subsection (or boxed remark) working out the classical fifth-force prediction
from Eq.~\eqref{eq:bianchi} for a representative test case (e.g., a classical test mass
near a superconducting chip), and quantify how it compares to known precision measurements.

(R3) Add a sentence to the falsification criterion noting that F2 ($\Delta p < 10^{-6}$\,Pa)
translates to $|\kteff|\cdot\Ment/R^2 < $ [specific number from the proposed experiment's
geometry], making the bound on $\kt$ dependent on $\asc$ and the Planck-scale
renormalization.

(R4) Either present the scalar-field action~\eqref{eq:action} as a minimal extension that
makes the conservation story covariant and natural, or add a ``model uncertainty'' box
noting that the conservation interpretation is unjustified without an action.

(R5) Concrete intermediate milestones for OP3: see Section~2.4.  The current statement
``derive $\asc$ from Lindblad'' is too vague for a publishable open problem.

\medskip\noindent\textbf{PUBLISHABILITY RECOMMENDATION}

Reject with major revisions required.  If C1--C5 are adequately addressed, the paper
would represent a publishable contribution as a theory proposal paper in Physical Review D
or Classical and Quantum Gravity.  It should not be submitted to Physical Review Letters
until OP3 is resolved (at least at the qualitative level in one platform).

\end{mdframed}

\subsection{Three-paper sequence}

Based on the referee critique, the most defensible publication strategy is a
\emph{three-paper sequence}, each making a single, independently defensible claim.

\medskip

\noindent\textbf{Paper~I: Theory Letter (4--6 pages, PRL/PRD Letters)}

\textit{Claim:} Dimensional consistency, $\kt$ estimate under explicit regulator hypothesis,
OP4 resolution, Bianchi-identity treatment with scalar-field extension, falsification criterion.

\textit{Target journal:} Physical Review D (Letters format) or Physical Review Research.

\textit{Prerequisite:} R1 (define $\Sent$), R4 (scalar-field action or explicit caveat),
R3 (reformulate falsification criterion).

\textit{What it does NOT need:} OP3 resolution, experimental result, any specific platform.

\medskip

\noindent\textbf{Paper~II: Technical Theory Backup (20--30 pages, PRD)}

\textit{Claim:} Full derivation of OP4, complete OP3 Lindblad model for one platform
(50-qubit transmon chain), regime hierarchy for $\asc$, OP8 Bianchi analysis with
quantitative bounds from equivalence-principle tests, OP10 renormalization group
power-counting argument.

\textit{Target journal:} Physical Review D (full article).

\textit{Prerequisite:} Numerical solution of Lindblad equation for cluster states,
$N \leq 20$ qubits.  Tensor network extension to $N = 50$.

\textit{Required new work:} Steps 1--2 of the OP3 milestone path (Section~2.4 above).

\medskip

\noindent\textbf{Paper~III: Experimental Proposal (10--15 pages, PRA or npj Quantum Info.)}

\textit{Claim:} A concrete 36-qubit neutral-atom Ramsey experiment that (a) sets a
bound on $|\kteff|\cdot\Ment/R^2$ and (b) is publishable as a platform-agnostic
experimental proposal even without performing the measurement.

\textit{Target journal:} Physical Review A (experimental) or npj Quantum Information.

\textit{Content:} Exact Hamiltonian and Lindblad parameters, Branch A/B state preparation
protocols, Ramsey measurement design, complete systematic error budget
(Table~\ref{tab:systematics}), and the null-result interpretation bound.

\textit{Required new work:} The mean-field energy matching calculation for cluster vs.\
product states.  The differential frequency shift from the local entanglement entropy density.

\medskip

\noindent\textbf{Timeline.}
\begin{center}
\begin{tabular}{@{}cll@{}}
\toprule
Months & Task & Paper \\
\midrule
1--3   & R1--R5 revisions to current white paper + scalar field action & I (theory letter) \\
3--6   & Lindblad numerics, $N \leq 20$, cluster state $\asc$ & II (regime hierarchy section) \\
6--12  & Tensor network extension to $N=50$; OP3 complete & II (full) \\
12--18 & Neutral atom experiment design + mean-field systematics & III (full) \\
\bottomrule
\end{tabular}
\end{center}

% ─────────────────────────────────────────────────────────────────────────────
\section*{Conclusion: What Is Solid, What Is Open, and What To Do Next}
\addcontentsline{toc}{section}{Conclusion}

\textbf{What is solid and should be published.}
The dimensional consistency of the modified field equation, the OP4 resolution
($\Delta a \propto \Ment/R^2$), the Bianchi identity consequences including
Eq.~\eqref{eq:bianchi}, the MICROSCOPE non-constraint argument, and the falsification
criterion structure (F1--F5 rephrased as a bound on $|\kteff|\Ment/R^2$).
These are internally consistent results that can be published as Paper~I now.

\textbf{What is open but not fatal.}
The $\kt = -1/4$ estimate.  This is motivated and internally consistent.  It cannot be
treated as a prediction until HRH-1 and HRH-2 are independently justified.  The paper
correctly labels it as a hypothesis.

\textbf{What is open and important to resolve.}
OP3 ($\asc$ for real platforms): the milestone path in Section~2.4 provides a concrete
plan requiring 3--6 months of numerical work.  OP10 (Planck-scale renormalization): harder,
requiring quantum gravity tools, but power-counting arguments in effective field theory may
provide qualitative guidance within 6--12 months.

\textbf{The single most important immediate action.}
Define $\Sent$ precisely as the internal inter-qubit entanglement entropy density, verify
that the cluster state maintains large $\asc$ under realistic decoherence for $N = 50$
(using Regime~II estimates), and restate the falsification criterion as a bound on the
observable combination $|\kteff|\Ment/R^2$.  These three changes transform the paper
from a work with significant ambiguities to one that is genuinely publishable.

% ─────────────────────────────────────────────────────────────────────────────
\bibliographystyle{unsrtnat}
\begin{thebibliography}{99}

\bibitem{Jacobson1995}
T.~Jacobson,
``Thermodynamics of spacetime: The Einstein equation of state,''
\textit{Phys.\ Rev.\ Lett.}\ \textbf{75}, 1260 (1995).

\bibitem{Verlinde2017}
E.~Verlinde,
``Emergent gravity and the dark universe,''
\textit{SciPost Phys.}\ \textbf{2}, 016 (2017).

\bibitem{Asenbaum2020}
P.~Asenbaum, C.~Overstreet, T.~Kim, J.~Curti, and M.~A.~Kasevich,
``Atom-interferometric test of the equivalence principle at the $10^{-12}$ level,''
\textit{Phys.\ Rev.\ Lett.}\ \textbf{125}, 191101 (2020).

\bibitem{MICROSCOPE2022}
P.~T.~Touboul \textit{et al.}\ (MICROSCOPE Collaboration),
\textit{Phys.\ Rev.\ Lett.}\ \textbf{129}, 121102 (2022).

\bibitem{Flammia2012}
S.~T.~Flammia, D.~Gross, Y.-K.~Liu, and J.~Eisert,
\textit{New J.\ Phys.}\ \textbf{14}, 095022 (2012).

\bibitem{Blume2017}
R.~Blume-Kohout \textit{et al.},
\textit{Nat.\ Commun.}\ \textbf{8}, 14485 (2017).

\bibitem{deHaro2001}
S.~de~Haro, S.~N.~Solodukhin, and K.~Skenderis,
``Holographic reconstruction of spacetime and renormalization in the AdS/CFT correspondence,''
\textit{Commun.\ Math.\ Phys.}\ \textbf{217}, 595 (2001).

\bibitem{Lukin2019}
H.~Bernien \textit{et al.},
``Probing many-body dynamics on a 51-atom quantum simulator,''
\textit{Nature}\ \textbf{551}, 579 (2017).

\bibitem{Saffman2016}
M.~Saffman,
``Quantum computing with atomic qubits and Rydberg interactions: Progress and challenges,''
\textit{J.\ Phys.\ B}\ \textbf{49}, 202001 (2016).

\bibitem{Braunstein1994}
S.~L.~Braunstein and C.~M.~Caves,
\textit{Phys.\ Rev.\ Lett.}\ \textbf{72}, 3439 (1994).

\end{thebibliography}

\end{document}
